\section{Background and Related Work}
\label{sec:related}

\paragraph{OOD detection for ANNs.}
The dominant families score a post-hoc statistic of a trained network: maximum softmax
probability~\cite{hendrycks2017msp}, temperature-scaled and input-perturbed
logits~\cite{liang2018odin}, free energy~\cite{liu2020energy}, rectified
activations~\cite{sun2021react}, virtual-logit matching~\cite{wang2022vim}, gradient
norms~\cite{huang2021gradnorm}, and distance/density on penultimate features
(Mahalanobis~\cite{lee2018mahalanobis}, deep $k$-NN~\cite{sun2022knn}). All presuppose
floating-point readouts, gradients, or repeated forward passes -- capabilities absent from
the spiking, feed-forward, integer-event execution model of neuromorphic hardware
(Section~\ref{sec:hardware} makes this precise).

\paragraph{OOD detection for SNNs (and how we differ).}
\advisor{point 1 -- confirm the task is new. Honest answer: OOD-for-SNN exists; OUR task framing is new.}
OOD detection \emph{inside} an SNN is not unprecedented. Most directly, Mart\'inez et al.~\cite{martinez2023snnood}
characterise the internal activations of hidden layers via \emph{spike-count} patterns and add
an explainability map, evaluated on image (including event-based) \emph{classification}
benchmarks. Our work is distinct on four axes that, in combination, define a new task:
(i)~we read the \emph{sub-threshold membrane potential}, not spike counts -- a continuous,
pre-threshold signal with strictly more information than the binarised spikes;
(ii)~the host model is an event-camera \emph{object detector} (Prophesee Gen1), not a
classifier, so there is no class posterior to lean on; (iii)~the distribution shifts are
\emph{event-stream corruptions} we define and release as a benchmark, not held-out classes;
and (iv)~we make the \emph{zero-compute / neuromorphic-register} feasibility argument central.
We therefore position the paper as introducing the \emph{event-camera membrane-statistics OOD}
task and the Gen1-C benchmark, while explicitly crediting prior SNN-OOD work.

\paragraph{Event-camera robustness and corruption benchmarks.}
Robustness benchmarks built from synthetic corruptions are standard for frame cameras
(ImageNet-C~\cite{hendrycks2019imagenetc}); the event-camera analogue is comparatively
underdeveloped. Gen1-C (Section~\ref{sec:corruptions}) fills this gap with corruptions
grounded in concrete sensor failure modes (Table~\ref{tab:corr-taxonomy}), seeded for exact
reproducibility, and operating on the standard histogram tensor so they compose with any
event model.

\paragraph{The host detector.}
Our membrane signal is read from the pretrained hybrid SNN--ANN detector of
Neftci et al.~\cite{neftci2025hybrid}, a strong Gen1 baseline that fuses a spiking front-end (SpikingJelly
PLIF~\cite{fang2021plif}) with an ANN detection head. We treat it as read-only and never
retrain it -- both for the OOD signal and for the motivation experiment of
Section~\ref{sec:motivation}.

\paragraph{Membrane potential as a hardware register.}
On Loihi~\cite{davies2018loihi}, BrainScaleS~\cite{schemmel2010brainscales}, and
SpiNNaker~\cite{furber2014spinnaker}, $V_{\mathrm{mem}}$ is state the neuron already
maintains; harvesting $\varphi$ from it is an $O(\text{channels})$ read. This is the basis of
the feasibility argument in Section~\ref{sec:hardware}.
